% Folder 78, batch 07 — archivists' pages 121–129 (end of the folder).
% Pass: Opus 5.5 (claude-opus-5-5), 2026-10-02 — first pass, unchecked against the pages by a human.
%
% Notes for maintainers: the variable of the polynomials F_n, S_n is written
% with a struck-through capital (read here as U); the cursive capital of the
% cyclotomic polynomial is rendered \Phi. Page 121 opens in the middle of a
% computation begun before the batch; page 129 changes subject (coverings and
% fundamental groups) and breaks off at the end of the folder.
\documentclass[11pt,a4paper]{article}
\input{../preamble/grothendieck}

\folder{78}
\batch{7}
\pages{121}{129}
\foldertitle{[Polygones réguliers et polynômes cyclotomiques] : notes manuscrites (s.d.).}
\dating{[à partir de 1977]}
\watermark{Édition de démonstration}

\begin{document}

\section{Transposée, produit tensoriel et valeurs propres}

\page{121}
\note{le calcul de cette page commence avant le début du lot : $u$, ${}^t u$, les bases $(e_i)$, $(e'_i)$ et l'élément $-2e'_1+e'_2$ sont introduits plus haut. La page est barrée de deux longs traits obliques parallèles ; elle est transcrite telle qu'elle se lit.}

\struck{$u$ \quad $u(e_1)$} \quad \struck{$\check u(x')\otimes$} \quad ${}^t u$

\struck{$u(x)\otimes x$}
\[
({}^t u\, x')\otimes x = x'\otimes u x
\]

\drawing{deux petites figures en haut à droite : un trapèze formé de deux droites horizontales et de deux obliques, puis le même trapèze avec des arcs marquant les angles et un point fortement marqué (une croix) à l'un des sommets inférieurs.}

\[
\bigl[{}^t u(-2e'_1+e'_2)\bigr]\otimes e_1 = (-2e'_1+e'_2)\otimes u e_1
\]
\[
{}^t u(e'_1) = \sum u_{1i}\, e'_i, \qquad {}^t u\, e'_2 = \sum u_{2i}\, e'_i, \qquad u e_1 = \sum_i u_{i1}\, e_i
\]

\struck{$\sum_i(-2u_{1i}+u_{2i})\,e'_i\otimes e_1$}
\[
= \sum(-2u_{i1})\, e'_1\otimes e_i + \sum u_{i1}\, e'_2\otimes e_i
\]

\struck{\ill{} $u_{1i}$} \quad \struck{\ill{} $u_{i1} =$}

\[
\begin{array}{cccc}
e'_1\otimes e_1 & e'_2\otimes e_1 & \cdots & e'_n\otimes e_1\\
e'_1\otimes e_2 & e'_2\otimes e_2 & & \\
\vdots & \vdots & & \\
e'_1\otimes e_n & e'_2\otimes e_n & &
\end{array}
\]

\[
\begin{cases} u_{i1} = 0 \\ -2u_{i1} = 0 \end{cases} \quad i\geqslant 2
\qquad\qquad
u_{2i} = 2u_{1i} \quad i\geqslant 3
\]
\[
-2u_{11}+u_{21} = -2u_{11}, \qquad -2u_{12}+u_{22} = u_{11}
\]
\uncertain{i.e.} $u_{21}=0$ déjà \emph{connu}

\[
\begin{array}{ccccc}
\lambda & u_{12} & u_{13} & \cdots & u_{1n}\\
0 & \lambda & 2u_{13} & \cdots & 2u_{1n}\\
\vdots & & & & \\
0 & & & &
\end{array}
\]
\note{au deuxième rang, l'entrée sous $u_{12}$ est surchargée (\uncertain{$\lambda$} écrit sur $u_{12}$) ; un ovale est tracé sous la matrice, vide.}

\marginal{écrit en long dans la marge gauche : $F\subset$ \ill{} ; \quad $a+\langle a,x'\rangle x$ ; \quad $x+\langle x,x'\rangle x$ \quad $-x$ ; plus bas, en travers, des symboles isolés ($x$, $\mathbb{Z}$, $x$) \ill{}.}

$\gamma$\uncertain{$!$}$(x-1)+$\ill{} $= 2\otimes 2$ \note{ligne au bas de la page, lecture très incertaine.}

\marginal{tête-bêche, au bas de la page à droite :}
\struck{$\lambda=1$}
\[
\lambda\,(T-1)^{\nu-1}(T+1)
\]
\[
(\lambda-1)^{\nu-1}(\lambda+1) = 0
\]
\uncertain{$\check E\otimes\check E$} \note{un $\lambda$ est écrit au-dessus, à droite.}

\section{Polynômes cyclotomiques et polynômes $F_n$}

\page{122}
$\Phi_n(T)$ polynôme minimal sur $\mathbb{Q}$ d'une racine primitive $n$-ième de $1$ \qquad $n\in\mathbb{N}^*$

$F_n(T)$ polynôme minimal de $\zeta+\zeta^{-1}$

Si $n\neq 1,2$, on a
\[
\Phi_n(T) = T^{\varphi'(n)}\, F_n(T+T^{-1}) \qquad (\varphi'(n)=\varphi(n)/2)
\]
\marginal{encadré, relié par une flèche : $\Phi_n, F_n\in\mathbb{Z}[T]$ pol.\ unitaires}

$\deg\Phi_n = \varphi(n)$ (indicateur d'Euler, $=\operatorname{card}((\mathbb{Z}/n\mathbb{Z})^*)$)
\[
\deg F_n = \begin{cases} 1 & \text{si } n=1,2 \quad (F_1(T)=T-2,\ F_2(T)=T+2)\\ \varphi(n)/2 & \text{si } n\geqslant 3\end{cases}
\]

\[
\prod_{\substack{d\in\mathbb{N}^*\\ d\mid n}} \Phi_d(T) = T^n-1
\]
donne mode de calcul par récurrence multiplicative sur $n$.

\[
\begin{array}{ll}
F_1(U) = U-2 & \Phi_1(T) = T-1\\
F_2(U) = U+2 & \Phi_2(T) = T+1\\
F_3(U) = U+1 & \Phi_3(T) = T^2+T+1\\
F_4(U) = U & \Phi_4(T) = T^2+1\\
F_5(U) = U^2+U-1 & \Phi_5(T) = T^4+T^3+T^2+T+1\\
F_6(U) = U-1 & \Phi_6(T) = T^2-T+1\\
F_7(U) = U^3+U^2-2U-1 & \Phi_7(T) = T^6+T^5+T^4+T^3+T^2+T+1\\
F_8(U) = U^2-2 & \Phi_8(T) = T^4+1\\
F_9(U) = U^3-3U+1 & \Phi_9(T) = T^6+T^3+1\\
F_{10}(U) = U^2-U-1 & \Phi_{10}(T) = T^4-T^3+T^2-T+1\\
\cdots & \cdots
\end{array}
\]
\note{la variable des $F_n$ et des $S_n$ est une capitale barrée, lue ici $U$ ; à la ligne de $F_4$, ce qui suit $U$ est noyé sous une rature.}

(donner aussi les cas $n=12,14,18\ldots$)

\[
\begin{cases}
\Phi_p(T) = 1+T+\cdots+T^{p-1} \quad \text{si } p \text{ premier}\\
F_p(T) = S_0(T)+S_1(T)+\cdots+S_{p'}(T)
\end{cases}
\]
où $p=2p'+1$ i.e.\ $p-1=2p'$ ($p$ est premier impair), et les $S_n(U)\in\mathbb{Z}[U]$ ($n\in\mathbb{N}^*$) sont caractérisés par
\[
\begin{cases} S_0(U) = 1\\ T^n+T^{-n} = S_n(T+T^{-1}) & n\geqslant 1\end{cases}
\]
donc
\[
\begin{cases} S_0(U)=1\\ S_1(U) = U\\ S_2(U) = U^2-2\\ S_3(U) = U^3-3U\\ \cdots\end{cases}
\]
Calcul récurrent des $S_i$ par formule \struck{$S_n(T)$}
\[
U^n = \sum_{0\leqslant i\leqslant[\frac n2]} \binom{n}{i} S_{n-2i}(U) \qquad n\geqslant 1
\]
d'où
\[
S_n(U) = U^n - \sum_{1\leqslant i\leqslant[\frac n2]}\binom{n}{i} S_{n-2i}(U)
\]
\note{pour $n$ pair, le terme $i=n/2$ devrait porter le coefficient $\frac12\binom{n}{n/2}$ ; la formule est transcrite telle qu'elle est écrite.}

\page{123}
\[
\begin{cases}
\Phi_{2p}(T) = 1-T+T^2+\cdots+(-1)^{p-1}T^{p-1}\\
F_{2p}(T) = (-1)^{p'}\,\Phi_p(-T)\\
\phantom{F_{2p}(T)} = \bigl(S_0(T)-S_1(T)+\cdots+(-1)^{p'}S_{p'}(T)\bigr)(-1)^{p'}
\end{cases}
\]
$p$ est premier impair
\note{au-dessus de la première ligne, relié à « $=1-T+\cdots$ » par un trait : « $=\Phi_p(-T)$ » suivi d'une rature, la lettre étant lue \uncertain{$F$} ou \uncertain{$\Phi$} ; dans la seconde ligne, le facteur $(-1)^{p'}$ est ajouté au-dessus, et la lettre devant $(-T)$ est la même capitale cursive, lue $\Phi_p$ quoique la formule appelle $F_p$.}

Sur un corps \add{alg.\ clos de} car.\ quelconque $p$, les racines de $F_n(T)$ sont les racines \emph{primitives} \struck{$n$-ièmes} \add{$n'$-ièmes} de $1$, chacune avec la même multiplicité, qui est
\[
\begin{cases}
\text{multiplicité } 1 & \text{si } (p,n)=1\\
\text{multiplicité } \varphi(p^r) & \text{si } n=p^r n',\ (n,n')=1,\ r\in\mathbb{N}
\end{cases}
\]
\note{interligne sous ces deux lignes : « primitives \struck{distinctes $n$-ièmes} \ldots{} racines \ldots{} $n'$-ièmes », partiellement barré, \ill{}. La condition $(n,n')=1$ est écrite ainsi ; on attend $(p,n')=1$.}

Si $(\zeta_i)$ \struck{\ill{}} sont les racines \uncertain{primitives $n'$-ièmes} de $1$, \ill{} l'on a
\[
\Phi_n(T) = \Bigl(\prod(T-\zeta_i)\Bigr)^{\varphi(p^r)} = \Phi_{n'}(T)^{\varphi(p^r)} \quad\text{dans } k[T]
\]
\struck{(\ill{}) \ill{} donc \ill{} de car.\ $p$ \ill{}}
\[
\Phi_n(T) = \Phi_{n'}(T)^{\varphi(p^r)} \quad\text{dans } \mathbb{F}_p(T)
\]
i.e.
\[
\Phi_n \equiv \Phi_{n'}^{\varphi(p^r)} \mod p
\]
En particulier
\[
\begin{cases}
\Phi_{p^r}(T) \equiv (T-1)^{\varphi(p^r)} \mod p\\
\Phi_p(T) \equiv (T-1)^{p-1} \mod p
\end{cases}
\]

Les racines de $F_n(T)$ dans $k$ sont les éléments de la forme $\zeta+\zeta^{-1}$ ($\zeta$ racine primitive $n'$-ième de $1$), chacune avec la multiplicité $\varphi(p^r)$ \struck{\ill{}} \add{\uncertain{sauf si} $n'=1$ ou $2$, où la multiplicité \uncertain{est} $\varphi(p^r)/2$} \uncertain{en particulier}
\[
\begin{cases}
F_n(T) \equiv F_{n'}(T)^{\varphi(p^r)} \mod p & \text{si } n'\neq 1,2\\
F_n(T) \equiv F_{n'}(T)^{\varphi(p^r)/2} \mod p & \text{si } n'=1 \text{ ou } 2
\end{cases}
\]
\[
\begin{cases}
F_{p^r}(T) \equiv \text{\struck{$F_p(T)^{\varphi(p)}$}}\ (T-2)^{\varphi(p^r)} \quad (p)\\
F_p(T) \equiv (T-2)
\end{cases}
\]
\note{la seconde ligne s'arrête là. L'exposant $\varphi(p^r)$ de la première contredit la règle $n'=1$ énoncée juste au-dessus ($\varphi(p^r)/2$) ; la page suivante reprend avec $\varphi(p^r)/2$.}

\page{124}
Multiplicité $\mu$ des racines de $F_n$ \add{$(n\geqslant 3)$} (dans $k$ de car.\ $p>0$),
\[
n = n'p^r \qquad ((n,n')=1,\ r\in\mathbb{N})
\]
1) $n'\neq 1,2$ (donc $F_{n'}(T)$ de degré $\varphi(n')/2$) : $\mu=\varphi(p^r)$, donc
\[
F_n(T) \equiv F_{n'}(T)^{\varphi(p^r)} \mod p
\]
2) $n'=1$ ou $2$ (donc $F_{n'}(T)$ de degré $1=\varphi(n')$) \struck{\ill{} $\neq 1,2$} et $p^r\neq 2$ i.e.\ $n\neq 4$ --- donc \struck{$\varphi(p^r)$} \add{$F_{p^r}(T)$} de degré $\varphi(p^r)/2$ \qquad $\Big|\ \mu = \varphi(p^r)/2$
\[
F_n(T) \equiv F_{n'}(T)^{\varphi(p^r)/2} \mod (p)
\quad
\begin{cases}
(T-2)^{\varphi(p^r)/2} & \text{si } n=p^r\\
(T+2)^{\varphi(p^r)/2} & \text{si } n=2p^r
\end{cases}
\]
\struck{3) $n=4=2^2$, $p=2$ \quad $F_4(T)\equiv(T-2)^2$ \ill{} $\equiv(T+2)^2\equiv T^2$ \ill{} $F_{n'=1}(T)^{\varphi(4)=4}$ \ill{} (\ill{} $1$)}
\note{cette ligne et la suivante sont annulées par un zigzag.}

En particulier, si $p$ premier $\neq 2$, ---.
\[
\begin{cases}
F_p(T) \equiv (T-1)^{p'} \mod p\\
F_{2p}(T) \equiv (T+1)^{p'} \mod p
\end{cases}
\qquad \Big|\ p' = (p-1)/2
\]
\note{sous $p'=(p-1)/2$, une courte formule raturée \ill{}. Écrit $T-1$ et $T+1$ ; la formule générale de la même page donne $T-2$ et $T+2$.}

\section{Les anneaux $A_p$, $B_p$ ; l'élément $\zeta-1$}

\page{125}
\[
\begin{array}{l}
\alpha\in B_{n(p)} = \mathbb{Z}[U]/F_p(U)\\
\zeta\in A_p = \mathbb{Z}[T]/\Phi_p(T)
\end{array}
\quad\Big|\ \text{normaux} \qquad \Big|\ p \text{ premier impair}
\]
\note{l'indice de $B$ est surchargé : \uncertain{$n(p)$} écrit sous un $p$.}

ils sont étales sur $\mathbb{Z}$ en les \struck{\ill{}} $\ell\in\operatorname{Spec}(\mathbb{Z})$, $\ell\neq p$ (car en car.\ $\ell$ les racines de $F_p$, $\Phi_p$ sont distinctes)

Ramifiés en $\ell=p$ (une seule racine en car.\ $p$)
\[
(\zeta-1)^{p-1} \equiv 0 \quad (p) \quad\text{dans } A_{(p)} = A_p
\]
\struck{\ill{} $\zeta^{p-1}-(p-1)\zeta^{p-2}+\frac{(p-1)(p-2)}{2}\zeta^{p-3}-\cdots 1$}

\struck{$1+T+\cdots+T^{p-1} \equiv (T-1)^{p-1}\quad(p)$}
\[
\begin{aligned}
(\zeta-1)^{p-1} &= (\zeta-1)^{p-1} - F_p(\zeta)\\
&= p\,\zeta^{p-2} + \Bigl(\tfrac{(p-1)(p-2)}{2}-1\Bigr)\zeta^{p-3}\\
&\qquad + \Bigl(-\tfrac{(p-1)(p-2)(p-3)}{6}-1\Bigr)\zeta^{p-4} - \cdots
\end{aligned}
\]
\note{il écrit $F_p(\zeta)$ là où le calcul porte sur $\Phi_p(\zeta)$.}

\drawing{un cercle, un rayon horizontal, et un court segment épais joignant l'extrémité du rayon à un point voisin du cercle.}
\[
\frac1p(\zeta-1)^{p-1} = \zeta^{p-2} + \cdots +
\]
\[
N\Bigl(\frac1p(\zeta-1)^{p-1}\Bigr) = \frac{1}{p^{p-1}}\, N(\zeta-1)^{p-1}.
\]
\[
N(\zeta-1) = \prod_1^{p-1}(\zeta_i-1) = \text{\struck{$\varepsilon$}}\prod(1-\zeta_i) = F_p(1) = \underbrace{1+\cdots+1}_{} = p
\]
\[
N\Bigl(\frac1p(\zeta-1)^{p-1}\Bigr) = 1
\]
\[
\boxed{(\zeta-1)^{p-1} = p\,u} \qquad u \text{ \emph{unité}}
\]
Donc $(\zeta-1)A \supset pA$ \quad \struck{$(\zeta-1)$}

d'où
\[
A/(\zeta-1)A = (\underbrace{A/pA}_{A_0})/(\zeta-1)A_0 = \bigl(\mathbb{F}_p[T]/(T-1)^{p-1}\bigr)/(T-1) = \mathbb{F}_p
\]
Donc $(\zeta-1)A$ est un idéal \emph{maximal}

\page{126}
$F_{p^r}$ \qquad $\varphi' = \varphi(p^r)/2$ \quad si $p^r\neq 2$
\note{le $2$ final est surchargé.}
\[
(\alpha-2)^{\varphi'} = p\,u
\]
\[
N(\alpha-2)^{\varphi'} = p^{\varphi'} N(u)
\]
\[
N(\alpha-2) = \prod(\alpha_i-2) = (-1)^{\varphi'}\prod(2-\alpha_i) = (-1)^{\varphi'} \underbrace{F_{p^r}(2)}_{p} = (-1)^{\varphi'}\,\Phi_{p^r}(1)
\]
\note{au-dessus de $\prod(2-\alpha_i)$, accroché par une accolade : $(\zeta'_i-1)(\zeta''_i-1)$ ; à droite de « $p$ » sous l'accolade : « \uncertain{$=(-1)^{\varphi'}$} ».}

$8+4-4-1$

\struck{$A_p B$} \quad $A_{p^r}$, $B_{p^r}$ sont normaux \uncertain{ou} Dedekind

$A_{n'p^r}$

$(\xi-1)$
\[
\mathbb{F}_p[T]/\Phi_n(T) = \mathbb{F}_p[T]/\text{\struck{$(T-1)^{\varphi(p^r)}$}}\ \Phi_{n'}(T)^{\varphi(p^r)}
\]
\[
\text{\struck{$=\mathbb{F}_p[T]/\Phi_{n'}(T^{\cdots})$}} \quad \prod_{(r,n')=1}(T-\xi_1^r)\Big/\prod(T-\zeta_i)^{\varphi(p^n)}
\]

$T-1$ \qquad \struck{$T^n$} \quad $1+T+\cdots+T^{n'-1}$ \qquad $T^{p^n}-\zeta = 0$

$n'$ premier à $p$

mais \uncertain{pas} \ill{} \uncertain{premier à} $p\ldots$ !

$T$ \quad $\zeta$ \quad $A_{n'}\subset A_n$ \qquad $\zeta$ !

$\zeta' = \zeta^{p^r}$ \qquad $\zeta-$\struck{$\zeta'$} \qquad \uncertain{$\zeta'^{(p^r)}$}
\[
N_{A_n/A_{n'}}(\zeta-1)
\]
\[
\underbrace{\mathbb{F}_p[T']/\Phi_{n'}(T')}\,[T]/(T^{p^n}-T') \xrightarrow{\ \text{cf.}\ } A_n\otimes\mathbb{F}_p
\]
\uncertain{Il y a des racines} \ill{} \uncertain{cycl.} $\Phi_n$. \quad \struck{\ill{}} $\Phi_n$ \quad $T^{\varphi(p^n)}+$
\note{l'étiquette de la flèche est lue \uncertain{cf.} ; sous la flèche, une suite de petits traits raturés.}

\page{127}
\[
(\alpha-2)^{p'} \equiv 0 \quad (p)
\]
\[
(\alpha-2)^{p'} = p\,u
\]
\[
N(\alpha-2)^{p'} = p^{p'} N(u)
\]
\[
N(\alpha-2) = \prod(\alpha_i-2) = (-1)^{p'}\prod(2-\alpha_i) = (-1)^{p'} F_p(2) = (-1)^{p'}\Phi_p(1) = (-1)^{p'}p
\]
\note{au-dessus de $\prod(2-\alpha_i)$ : $(\zeta'_i-1)(\zeta''_i-1)$ ; en dessous, deux essais raturés \struck{$\zeta'_i\zeta''_i$}, \struck{(\ill{})}. La dernière égalité est écrite $=(-1)^{p'}p$ sous la précédente.}

\struck{$\alpha_i-2 = \zeta_i^2+\zeta_i^{-2}-2 = (\zeta_i-\zeta_i^{-1})^2$}

$F_5$ \quad $8+4-4-1$ \qquad $\boxed{F_p(2)=p}$ ?
\[
(\zeta-1)(\zeta^{-1}-1) = 1+1-\alpha
\]
\struck{$A = B[T]$}
\[
B[T]/(T^2-\alpha T+1) \xrightarrow{\ \text{iso ?}\ } A
\]

\[
(\zeta-1)^{\varphi} \equiv 0 \quad (p)
\]
\note{un « $2$ » est écrit au-dessus du $\zeta$ ; devant le $0$, un signe raturé \ill{}.}
\[
(\xi-1)^{\varphi} = p\,u
\]
\[
N(\xi-1)^{\varphi} = p^{\varphi} N(u)
\]
\[
\prod(\zeta_i-1) = (-1)^{\varphi}\prod_i(1-\zeta_i) = (-1)^{\varphi}\,\underbrace{\Phi_{p^r}(1)}_{p}
\]
\[
\begin{cases}
\displaystyle\prod_{\substack{d\mid p^r\\ d\neq 1}} \Phi_d(T) = \text{\struck{$T^n$}}\ 1+T+\cdots+T^{p^r-1} = p^r \quad p\\
\text{\struck{$\displaystyle\prod_{d\mid p^r} p = \Phi_p(1)\Phi_{p^2}(1)\cdots$}} \quad p\,p^?-p^r\\
\Phi_{p^r}(1) = p
\end{cases}
\]
\note{la ligne du milieu est barrée sauf sa fin, lue \uncertain{$p\,p^?-p^r$} ; l'égalité $=p^r$ de la première ligne vaut pour $T=1$.}

\section{Les anneaux $A_n$, $B_n$ et l'involution $\sigma$}

\page{128}
\[
\xi\in A_{n'n''=n} \supset A_{n''}\ni\zeta'' \qquad ((n',n'')=1)
\]
\[
\cup \qquad \zeta'\in A_{n'}
\]
\[
G = (\mathbb{Z}/n\mathbb{Z})^* \simeq G'\times G'', \qquad G'\simeq(\mathbb{Z}/n'\mathbb{Z})^*,\quad G''\simeq(\mathbb{Z}/n''\mathbb{Z})^*
\]
\[
\underbrace{A_{n'}\otimes_{\mathbb{Z}} A_{n''}} \xrightarrow{\ \sim\ } A_{n'n''}
\]
normal par hyp.\ de réc.

$B$ \qquad $\varphi(1)=1$, \quad $\varphi(2)=1$
\[
\zeta_i\eta_j + \zeta_i^{-1}\eta_j^{-1}
\]
\[
B_n[T]/(T^2-\alpha T+1) \longrightarrow A_n \qquad
\begin{cases} \alpha_n\longmapsto \xi_n+\xi_n^{-1}\\ T\longmapsto \xi_n\end{cases}
\]
\[
\Big|\quad
\begin{aligned}
&\mu_n\supset\mu_n^* = \operatorname{Spec}\mathbb{Z}[T]/\Phi_n(T)\\
&\mu_{nn'}^*\simeq\mu_n^*\times\mu_{n'}^*\\
&\mu_n^*/\{\pm1\} = \tilde\mu_n^*
\end{aligned}
\]
\note{dans la dernière ligne, un symbole raturé précède $\tilde\mu_n^*$.}
\[
B_n = A_n^{\sigma}
\]
\note{l'exposant de $A_n$ est un signe peu net, lu \uncertain{$\sigma$}.}
\[
\begin{aligned}
\sigma(F+GT) &= F+G(U-T)\\
&= (F+GU)-GT
\end{aligned}
\]
\[
\sigma(F+GT) = F+GT \iff
\begin{cases} GU=0 & G=-G\\ \iff G=0 \end{cases}
\]

\begin{tikzcd}[column sep=small, row sep=small]
A=\mathbb{Z}[T,T^{-1}] \arrow[r] & \mathbb{Z}[T]/\Phi_n(T)=A_n\\
B=\mathbb{Z}[U]=\mathbb{Z}[T]^{\sigma} \arrow[r] \arrow[u, hook] & \mathbb{Z}[U]/F_n(U)=B_n \arrow[u, hook]
\end{tikzcd}

\[
\begin{array}{l}
U = T+T^{-1}\\
\sigma(T) = T^{-1}
\end{array}
\ \Big|\
F_n(U) = F_n(T+T^{-1}) = T^{-\varphi'}\Phi_n(T)
\]

\begin{tikzcd}[column sep=small, row sep=small]
\mathbb{G}_m \arrow[d] & \supset\ \mu_n\ \supset\ \mu_n^* & \\
\mathbb{E}^1 & \tilde\mu_n/\{\pm1\} \arrow[l] & \tilde\mu_n^* \arrow[l]
\end{tikzcd}

$\mathbb{E}^1\simeq\mathbb{G}_m/\sigma$ \qquad $T' = U-T$
\note{le $\mathbb{E}$ est une capitale à double barre ; lecture \uncertain{$\mathbb{E}^1$} ; le $T$ de $T'=U-T$ est souligné.}

\emph{$n\neq 1,2$}
\[
A \xleftarrow{\ \sim\ } B[S]/(S^2-US+1) \qquad\text{\struck{$\mathbb{Z}[T]/(T^n-1)$}}
\]
i.e.\ $A$ \uncertain{est libre} sur $B$ de base $1,T$
\[
= B[\dot S] = B.1\oplus B.\dot S
\]
\uncertain{$T^{2n}=T^n$}
\[
\text{\struck{$F+T$}}\quad F(U)+T\,G(U) = 0 \ \Longrightarrow\ F(U)+T^{-1}G(U) = 0
\]
\[
T\,G(U) = T^{-1}G(U), \qquad (T-T^{-1})\,G(U) = 0 \qquad G(U) = 0,
\]
\ill{}

\section{Revêtements ramifiés et groupes fondamentaux}

\page{129}
\note{changement de sujet : la page porte sur des revêtements d'une base $S$ ramifiés en deux points, et sur les groupes fondamentaux correspondants. Elle est la dernière du dossier ; le raisonnement y reste en suspens.}

\begin{tikzcd}[column sep=small, row sep=small]
B \arrow[dr, no head] & & B' \arrow[dl, no head]\\
& A &
\end{tikzcd}

\begin{tikzcd}[column sep=small, row sep=small]
& X\times_S X \arrow[dl, no head] \arrow[dr, no head] & \\
X \arrow[dr, no head] & & X' \arrow[dl, no head]\\
& S &
\end{tikzcd}

$G$ \quad $G'$ \qquad $s,s'\in S$

$X$ tot.\ ramifié en $s$, étale en dehors

$X'$ tot.\ ramifié en $s'$, étale en dehors

$X$, $X'$ \uncertain{normaux} irréd.

$X\times_S X'$ \uncertain{est connexe} ?

\begin{tikzcd}[column sep=small, row sep=small]
& \pi_1(S-\{s,s'\}) \arrow[dl] \arrow[dr] & \\
\pi_1(S-s) \arrow[dr] & & \pi_1(S-s') \arrow[dl]\\
& \pi_1(S) &
\end{tikzcd}

\begin{tikzcd}[column sep=small, row sep=small]
& G*\mathbb{Z}*\mathbb{Z} \arrow[dl] \arrow[dr] & \\
X*\mathbb{Z} \arrow[dr] & & G*\mathbb{Z} \arrow[dl]\\
& G &
\end{tikzcd}

\note{le sommet de gauche est lu \uncertain{$X*\mathbb{Z}$} ; on attendrait $G*\mathbb{Z}$.}

$G*G'$

\drawing{un petit carré raturé : $X_s$, $X'_s$ et des flèches vers $S$, couvert de traits de biffure.}

\drawing{deux arcs se terminant chacun par un point épais, au-dessus d'un segment horizontal marqué à droite d'un trait vertical et de la lettre $H$ (un astérisque raturé à gauche).}

$I\subset H_s\subset G\times G'$ \qquad surj.\ vers $G$, surj.\ vers $G'$

$Z\subset X\times_S X$ \qquad $Z\to X$ : \ill{} \uncertain{revêt.}\ tot.\ ramifié en les pts de $X_s$ \qquad $Z\to X'$ : tot.\ ram.\ en les pts de $X'_s$
\note{le second facteur est écrit $X$ ; la flèche va vers $X'$.}

$(\mathbb{Z}/4\mathbb{Z})^*$ d'ordre $2$ \qquad $\mathbb{Z}/$

\end{document}
